
\documentclass{article}
\usepackage{colortbl}
\usepackage{makecell}
\usepackage{multirow}
\usepackage{supertabular}

\begin{document}

\newcounter{utterance}

\centering \large Interaction Transcript for game `cladder', experiment `full\_v1.5\_default', episode 948 with Qwen3.6{-}27B.
\vspace{24pt}

{ \footnotesize  \setcounter{utterance}{1}
\setlength{\tabcolsep}{0pt}
\begin{supertabular}{c@{$\;$}|p{.15\linewidth}@{}p{.15\linewidth}p{.15\linewidth}p{.15\linewidth}p{.15\linewidth}p{.15\linewidth}}
    \# & \multicolumn{2}{c}{Player} && \multicolumn{2}{c}{Game Master} \\
    \hline

    \theutterance \stepcounter{utterance}  
    & & & \multicolumn{4}{p{0.6\linewidth}}{
        \cellcolor[rgb]{0.9,0.9,0.9}{
            \makecell[{{p{\linewidth}}}]{
                \texttt{\tiny{[P1$\langle$GM]}}
                \texttt{You are an expert in causal inference. The following question is not a typical commonsense query, but rather a meticulously designed question created by a professor specializing in causal inference, intended to assess the students' mastery of the course content.} \\
\texttt{Imagine a self{-}contained, hypothetical world with only the following conditions, and without any unmentioned factors or causal relationships: Obesity has a direct effect on diabetes and lifespan. Smoking has a direct effect on diabetes and lifespan. Diabetes has a direct effect on lifespan. Smoking is unobserved. For normal weight people and without diabetes, the probability of long lifespan is 91\%. For normal weight people and with diabetes, the probability of long lifespan is 71\%. For obese people and without diabetes, the probability of long lifespan is 69\%. For obese people and with diabetes, the probability of long lifespan is 43\%. For normal weight people and nonsmokers, the probability of having diabetes is 28\%. For normal weight people and smokers, the probability of having diabetes is 20\%. For obese people and nonsmokers, the probability of having diabetes is 57\%. For obese people and smokers, the probability of having diabetes is 70\%. The overall probability of smoker is 8\%. If we disregard the mediation effect through diabetes, would obesity negatively affect lifespan?} \\
\texttt{Start your answer with 'Yes' or 'No', followed by additional reasoning or evidence to support your explanation.} \\
            }
        }
    }
     \\ \\

    \theutterance \stepcounter{utterance}  
    & \multicolumn{4}{p{0.6\linewidth}}{
        \cellcolor[rgb]{0.561,0.737,0.561}{
            \makecell[{{p{\linewidth}}}]{
                \texttt{\tiny{[P1$\rangle$GM]}}
                \texttt{No} \\
\\ 
\texttt{To determine whether obesity negatively affects lifespan when disregarding the mediation effect through diabetes, we must isolate the **direct effect** of obesity on lifespan. In causal inference terms, this corresponds to the Controlled Direct Effect (CDE) or the natural direct effect, where we hold the mediator (diabetes) constant and compare the outcome (lifespan) between the treatment (obese) and control (normal weight) groups.} \\
\\ 
\texttt{The problem provides the probability of a long lifespan conditional on weight status and diabetes status. Since we are disregarding the mediation effect, we compare the probabilities of long lifespan for obese individuals versus normal weight individuals **within the same diabetes status**.} \\
\\ 
\texttt{**1. Compare outcomes for individuals WITHOUT diabetes:**} \\
\texttt{*   Probability of long lifespan for Normal Weight, No Diabetes: \$P(L=1 | W=0, D=0) = 0.91\$} \\
\texttt{*   Probability of long lifespan for Obese, No Diabetes: \$P(L=1 | W=1, D=0) = 0.69\$} \\
\\ 
\texttt{Difference (Obese {-} Normal Weight): \$0.69 {-} 0.91 = {-}0.22\$} \\
\texttt{In this subgroup, obesity is associated with a **22 percentage point decrease** in the probability of a long lifespan.} \\
\\ 
\texttt{**2. Compare outcomes for individuals WITH diabetes:**} \\
\texttt{*   Probability of long lifespan for Normal Weight, Diabetes: \$P(L=1 | W=} \\
            }
        }
    }
    & & \\ \\

    \theutterance \stepcounter{utterance}  
    & & & \multicolumn{2}{p{0.3\linewidth}}{
        \cellcolor[rgb]{0.9,0.9,0.9}{
            \makecell[{{p{\linewidth}}}]{
                \texttt{\tiny{[GM$|$GM]}}
                \texttt{No} \\
            }
        }
    }
    & & \\ \\

    \theutterance \stepcounter{utterance}  
    & & & \multicolumn{2}{p{0.3\linewidth}}{
        \cellcolor[rgb]{0.9,0.9,0.9}{
            \makecell[{{p{\linewidth}}}]{
                \texttt{\tiny{[GM$|$GM]}}
                \texttt{yes} \\
            }
        }
    }
    & & \\ \\

    \theutterance \stepcounter{utterance}  
    & & & \multicolumn{2}{p{0.3\linewidth}}{
        \cellcolor[rgb]{0.9,0.9,0.9}{
            \makecell[{{p{\linewidth}}}]{
                \texttt{\tiny{[GM$|$GM]}}
                \texttt{game\_result = LOSE} \\
            }
        }
    }
    & & \\ \\

\end{supertabular}
}

\end{document}
